% Five more instruments across the five books. Numbers from build/compare/numbers.tex.
\chapter{Five more ways of looking}
\label{ch:more}

The walk turns on feeling. A book has other currents running through it: its
breathing, its questions, the people it speaks of, the hours it keeps, the
light it falls into. Each of the five instruments in this chapter follows one
of them, and each is drawn for all five books.

\clearpage
\section{Breath}

\stacked{breath}{the-trial}
\stacked{breath}{bartleby}
\stacked{breath}{yellow-wallpaper}
\stacked{breath}{heart-of-darkness}
\stacked{breath}{pride-and-prejudice}

\noindent Read aloud, a book makes you breathe where its punctuation lets you.
Here every clause (the words between one comma, semicolon, dash or full stop
and the next) is a vertical stroke as tall as the clause is long, and the
book runs from left to right in two rows. A spiky skyline is a book that holds
its breath.

\bk{the-trial} breathes longest. Its clauses average \lab{c-the-trial-clause}
words, and \lab{c-the-trial-longbreath}\,\% of them run to twenty words or
more, against \lab{c-pride-and-prejudice-longbreath}\,\% in
\bk{pride-and-prejudice} and \lab{c-bartleby-longbreath}\,\% in \bk{bartleby}.
Kafka's sentences unroll in long, unbroken reasoning that rarely pauses
(``K. waited a little while, looked from his pillow at the old woman who lived
opposite and who was watching him\ldots''). Austen, Melville and Conrad punctuate
in shorter breaths.

\clearpage
\section{The punctuation sky}

\stacked{sky}{the-trial}
\stacked{sky}{bartleby}
\stacked{sky}{yellow-wallpaper}
\stacked{sky}{heart-of-darkness}
\stacked{sky}{pride-and-prejudice}

\noindent Every sentence is one star in this sky, read like a page from left to
right and top to bottom: a small dot if it ends in a full stop, an open ring if
it asks a question, an upright bar if it exclaims. A sky is as large as its book.

\bk{yellow-wallpaper} is full of bars: \lab{c-yellow-wallpaper-exclaim}\,\% of
its sentences end in an exclamation mark, ten times the rate of
\bk{the-trial} (\lab{c-the-trial-exclaim}\,\%). The narrator cannot raise her
voice to anyone, so she raises it on the page. \bk{bartleby} asks most
(\lab{c-bartleby-questions}\,\%), and its densest constellation of rings,
\lab{c-bartleby-qrun} questions in twenty sentences, is the lawyer's helpless
interrogation of his clerk: ``What earthly right have you to stay here? Do you
pay any rent?'' In \bk{the-trial} the rings cluster most tightly in
chapter~\lab{c-the-trial-qrun-ch}, \lab{c-the-trial-qrun} questions in twenty
sentences. No one else asks them; K. is alone, asking himself: ``What would
his life be like in the days ahead? \ldots{} Would he survive that?'' In
\bk{pride-and-prejudice} the tightest cluster is chapter~\lab{c-pride-and-prejudice-qrun-ch},
the family questioning each other after Lydia's elopement.

\clearpage
\section{The pronoun walk}

\begin{center}
\fiverow{pronouns}{22mm}
\end{center}

\noindent A second walk, steered not by feeling but by person. Every pronoun is
one step: \emph{I} and \emph{me} step north, \emph{you} south, \emph{he} east,
\emph{she} west, \emph{we} north-east and \emph{they} south-west. Each book
starts at the black square. Where a walk ends up says whose book it is.

\begin{table}[h]
\centering\figfont\footnotesize
\setlength\tabcolsep{4pt}
\begin{tabular}{@{}lrrrrr@{}}
\toprule
 & \emph{Trial} & \emph{Bartleby} & \emph{Wallpaper} & \emph{Darkness} & \emph{Pride}\\
\midrule
\emph{I} (\%) & \lab{c-the-trial-pr-I} & \lab{c-bartleby-pr-I} & \lab{c-yellow-wallpaper-pr-I} & \lab{c-heart-of-darkness-pr-I} & \lab{c-pride-and-prejudice-pr-I}\\
\emph{you} & \lab{c-the-trial-pr-you} & \lab{c-bartleby-pr-you} & \lab{c-yellow-wallpaper-pr-you} & \lab{c-heart-of-darkness-pr-you} & \lab{c-pride-and-prejudice-pr-you}\\
\emph{he} & \lab{c-the-trial-pr-he} & \lab{c-bartleby-pr-he} & \lab{c-yellow-wallpaper-pr-he} & \lab{c-heart-of-darkness-pr-he} & \lab{c-pride-and-prejudice-pr-he}\\
\emph{she} & \lab{c-the-trial-pr-she} & \lab{c-bartleby-pr-she} & \lab{c-yellow-wallpaper-pr-she} & \lab{c-heart-of-darkness-pr-she} & \lab{c-pride-and-prejudice-pr-she}\\
\emph{we} & \lab{c-the-trial-pr-we} & \lab{c-bartleby-pr-we} & \lab{c-yellow-wallpaper-pr-we} & \lab{c-heart-of-darkness-pr-we} & \lab{c-pride-and-prejudice-pr-we}\\
\emph{they} & \lab{c-the-trial-pr-they} & \lab{c-bartleby-pr-they} & \lab{c-yellow-wallpaper-pr-they} & \lab{c-heart-of-darkness-pr-they} & \lab{c-pride-and-prejudice-pr-they}\\
Heading & \lab{c-the-trial-heading} & \lab{c-bartleby-heading} & \lab{c-yellow-wallpaper-heading} & \lab{c-heart-of-darkness-heading} & \lab{c-pride-and-prejudice-heading}\\
\bottomrule
\end{tabular}
\end{table}

\noindent \bk{the-trial} walks due east: nearly half its pronouns
(\lab{c-the-trial-pr-he}\,\%) are \emph{he}. It is told about K., never by him.
\bk{yellow-wallpaper} walks due north, \emph{I} after \emph{I}
(\lab{c-yellow-wallpaper-pr-I}\,\%). \bk{bartleby} and \bk{heart-of-darkness}
both head north-east, between the \emph{I} of a narrator and the \emph{he} of
the man he cannot understand, Bartleby or Kurtz. Only
\bk{pride-and-prejudice} turns west: it is the one book of the five in which
\emph{she} is the commonest pronoun.

\clearpage
\section{The clock}

\begin{center}
\fiverow{clock}{22mm}
\end{center}

\noindent Every mention of a time of day (dawn, morning, breakfast, noon,
afternoon, dinner, evening, dusk, supper, night, midnight) is a dot on a
24-hour dial, with midnight at the top and noon at the bottom. Its distance from
the centre is how far into the book it comes, so the first chapters sit near the
middle and the last near the rim.

\bk{bartleby} keeps office hours. Its times average \lab{c-bartleby-hour}, its
commonest are \lab{c-bartleby-time-top}, and only \lab{c-bartleby-night}\,\% of
them fall at night. \bk{yellow-wallpaper} lives at night:
\lab{c-yellow-wallpaper-night}\,\% of its times are night-time, averaging
\lab{c-yellow-wallpaper-hour}, the hours when the narrator lies awake and the
pattern moves. \bk{heart-of-darkness} begins at dusk on the Thames and stays
there (\lab{c-heart-of-darkness-night}\,\% night, averaging
\lab{c-heart-of-darkness-hour}). \bk{the-trial} and \bk{pride-and-prejudice}
both keep to morning and evening: arrests, calls and dinners. But \bk{the-trial}
has more of the night (\lab{c-the-trial-night}\,\% against
\lab{c-pride-and-prejudice-night}\,\%): the cathedral darkening around K.,
and the evening he is fetched.

\clearpage
\section{The light meter}

\noindent{\figfont\scriptsize\itshape The Trial}\\[1.5mm]
\includegraphics[width=\textwidth]{\cdir light-the-trial.png}\par\medskip
\noindent{\figfont\scriptsize\itshape Bartleby, the Scrivener}\\[1.5mm]
\includegraphics[width=\textwidth]{\cdir light-bartleby.png}\par\medskip
\noindent{\figfont\scriptsize\itshape The Yellow Wallpaper}\\[1.5mm]
\includegraphics[width=\textwidth]{\cdir light-yellow-wallpaper.png}\par\medskip
\noindent{\figfont\scriptsize\itshape Heart of Darkness}\\[1.5mm]
\includegraphics[width=\textwidth]{\cdir light-heart-of-darkness.png}\par\medskip
\noindent{\figfont\scriptsize\itshape Pride and Prejudice}\\[1.5mm]
\includegraphics[width=\textwidth]{\cdir light-pride-and-prejudice.png}\par\bigskip

\noindent Some words are light (\emph{light, sun, bright, white, lamp, candle,
shining}) and some are dark (\emph{dark, shadow, black, night, gloom, fog,
mist}). Each band reads through a book from left to right and prints the balance
between them in the stretch around each point: white where the book is bright,
black where it is dark, in the same fine dots used for every grey in this book.
Where there are too few such words to judge, the band is left blank.

\bk{heart-of-darkness} uses far more of these words than any other book
(\lab{c-heart-of-darkness-light} light, \lab{c-heart-of-darkness-dark} dark),
and its band is the most strongly marked. Its darkest stretch comes
\lab{c-heart-of-darkness-darkest-at}\,\% of the way through, in the second
part, as the steamer pushes upriver towards Kurtz. The surprise is
\bk{the-trial}, which on balance is a lighter book than its reputation
(\lab{c-the-trial-light} light words against \lab{c-the-trial-dark} dark).
Its light is mostly carried: candles and lamps held up in dim offices, attics
and the cathedral. \bk{pride-and-prejudice} has so few of either
(\lab{c-pride-and-prejudice-light} and \lab{c-pride-and-prejudice-dark}) that
most of its band is blank. Austen's world is lit evenly and does not remark
on it.
